\documentclass[12 pt, letterpaper]{hmcpset}
\usepackage{hw-preamble}

\name{Jayden Thadani}
\class{MATH 139 --- Fourier Analysis}
\assignment{Homework 9}
\duedate{4th April 2025}

\begin{document}

\begin{problem}[S \& S \S 5.5 \#15.]
	This exercise provides another example of periodisation.
	\begin{enumerate}
		\item Apply the Poisson summation formula to the function $g$ in Exercise 2 to obtain
			\[
				\sum_{n = -\infty}^{\infty} \frac{1}{(n + \alpha)^{2}} = \frac{\pi^{2}}{\sin^{2}{(\pi \alpha)}}
			\]
			whenever $\alpha$ is real, but not equal to an integer.

		\item Prove as a consequence that
			\[
				\sum_{n = -\infty}^{\infty} \frac{1}{n + \alpha} = \frac{\pi}{\tan{(\pi \alpha)}}
			\]
			whenever $\alpha$ is real but not equal to an integer.
	\end{enumerate}
\end{problem}

\pagebreak

\begin{problem}[S \& S \S 5.5 \#17.]
	The \emph{gamma function} is defined for $s > 0$ by
	\[
		\Gamma(s) = \int_{0}^{\infty} e^{-x} x^{s - 1} \, dx.
	\]
	\begin{enumerate}
		\item Show that for $s > 0$ the above integral makes sense, that is, that the following two limits exist:
			\[
				\lim_{\delta \to 0^{+}} \int_{\delta}^{1} e^{-x} x^{s - 1} \, dx \quad \text{and} \quad \lim_{A \to \infty}  \int_{1}^{A} e^{-x} e^{s - 1} \, dx.
			\]

		\item Prove that $\Gamma(s + 1) = s \Gamma(s)$ whenever $s > 0$, and conclude that for every integer $n \ge 1$ we have $\Gamma(n + 1) = n!$.

		\item Show that
			\[
				\Gamma(1 / 2) = \sqrt{\pi} \quad \text{and} \quad \Gamma(3 / 2) = \frac{\sqrt{\pi}}{2}.
			\]
	\end{enumerate}
\end{problem}

\pagebreak

\begin{problem}[S \& S \S 5.5 \#21.]
	Suppose that $f$ is continuous on $\mathbb{R}$. Show that $f$ and $\hat{f}$ cannot both be compactly supported unless $f = 0$. This can be viewed in the same spirit as the uncertainty principle.
\end{problem}

\end{document}
